\documentclass[book.tex]{subfiles}

\begin{document}

\chapter{Nonequilibrium Statistical Mechanics and Differentiable Phonon Simulations}
\label{chap:lattice_graphconvs}
\noindent\rule{11cm}{0.4pt} \\
\textbf{Prerequisites:} \autoref{chap:graphconvs}, \autoref{chap:molecular_ml}   \\
\textbf{Difficulty Level:} *** \\
\noindent\rule{11cm}{0.4pt}
\vspace{5mm}

Nonequilibrium statistical mechanical systems require new mathematical tools to model and simulate. In this chapter, we will learn the basics of modeling nonequilibrium systems with differentiable systems, in particular by studying the test problem of modeling nanoscale heat transport.





%"""Nanostructures are ideal candidates for heat management and conversion thanks to their ability of tuning thermal transport via phonon-boundary scattering. Modeling nanoscale heat transport, however, is challenging because Fourier theory breaks down [1]. In this talk, we report on a methodology for large-scale optimization of non-Fourier thermal transport in nanostructures, based upon the adjoint phonon Boltzmann transport equation and density-based topology optimization [2]. To attain arbitrary material distributions, we develop the transmission interpolation model, an interface-based method that allows for smooth interpolation between void and solid regions. We apply the method to two exemplary cases: i) Tailoring the effective thermal conductivity of a periodic system and ii) maximizing phonon size effects. Special emphasis will be given to numerical challenges and implementation details. The code for this work will be made available in the next release of OpenBTE [3]."""

\section{Macroscale Heat Conduction}
Fourier's law of heat conduction dictates that the rate of heat transfer at a point is proportional to the negative gradient in the temperature at the point. Mathematically we can describe this law as

\begin{align}
    q &= -\kappa \nabla T
\end{align}

where $\kappa$ is a quantity known as the material's conductivity. $\kappa$ is usually a scalar but can be a tensor for non-uniform methods. 

Fourier's law for heat transport is valid across a broad range of systems but breaks down at the nanoscale. The major challenge is that Fourier's law doesn't take phonons, the main heat carrier at the nanoscale into account. The phonon Boltzmann transport equation more accurately describes the physics of the system. 

%Optimizing thermal conductivity could prove important for optimizing thermoelectric systems. Being able to design nanostructures could enable breakthrough results but 
Modeling heat conduction using the phonon Boltzmann transport equation is very computationally challenging. Differentiable physics techniques can be used to construct a fully differentiable model of the phonon Boltzmann transport equation by constructing a differentiable representation of the density
\cite{romano2022differentiable}
\cite{romano2021openbte}.


\section{Boltzmann Transport Equation}

Suppose we have a distribution over space

\begin{align}
f(\vec{r}, \vec{v}, t)
\end{align}
where $\vec{r} \in \mathbb{R}^3, \vec{v} \in \mathbb{R}^3, t \in \mathbb{R}$

At equilibrium, this distribution is given by the Boltzmann equation
\begin{align}
    f(E) &= \frac{1}{\exp{\frac{E-E_F}{k_B T}} + 1}
\end{align}
The Boltzmann equation allows for modeling the direct form of the density in non-equilibrium situations as
\begin{align}
    \frac{df}{dt} &= \left ( \frac{df}{dt} \right )_{\mathrm{force}} + \left ( \frac{df}{dt} \right )_{\mathrm{diff}} + \left ( \frac{df}{dt} \right )_{\mathrm{coll}}
\end{align}
where the three terms on the right correspond to evolution due to external forces, diffusion, and particle collisions respectively. The form this commonly takes is
\begin{align}
    \frac{df}{dt} + \frac{p}{m} \cdot \nabla f + F \cdot \frac{df}{dp} &= \left ( \frac{df}{dt} \right )_{\mathrm{coll}}
\end{align}

The phonon Boltzmann transport equation applies the general theory of Boltzmann transport equations to phonons. We suppose we have a distribution
\begin{align}
    f_{qp}(r, t)
\end{align}
for the phonon state $(q, p)$ where $q$ is the wave number and $p$ is the polarization. 

The temperature gradient $\nabla T$ induces a heat current $J$ in the system. $J$ can then be obtained by the equation
\begin{align}
    J &= \frac{1}{(2\pi)^3}\sum_p \int f_{qp} h \omega_{qp} v_{qp} dq
\end{align}
Here $\omega$ is the angular frequency. At a system steady state, we have the stability condition
\begin{align}
    \frac{df_{qp}}{dt} &= 0
\end{align}
This term is commonly broken down into two terms
\begin{align}
    \frac{df_{qp}}{dt} &=  \frac{df_{qp}}{dt} \vert_{\mathrm{diffusion}} + \frac{df_{qp}}{dt} \vert_{\mathrm{scattering}}
\end{align}
The diffusion term is given by the equation
\begin{align}
    -v_p(q) \cdot \nabla T \frac{\partial f_{qp}}{\partial T}
\end{align}
The scattering term is harder to understand and is formulated in terms of perturbation theory.

\section{The Basics of Nanoscale Heat Transport}

The phonon Boltzmann transport equation is often used to model heat transport at the nanoscale. Boltzmann transport equations more generally are used to model the behavior of systems not in equilibrium. 

\begin{align}
    -C_p  v_p \cdot \nabla \Delta T_p(r) &= \sum_q W_{pq} \Delta T_q(r)
\end{align}

Here, $p$ and $q$ are the phonon branch and wave vector. $\Delta T_q$ is the phonon pseudo-temperature. $C_p$ is the mode specific heat capacity and $v_p$ is the group velocity. $W_{p q}$ is a representation of the scattering operator. This equation can be solved by a finite element style method.

We start by assuming that a bulk material is characterized by its thermal conductivity tensor $\kappa$ and its phonon mean free path distribution $\Lambda$. The mean free path for phonons is a measure of the material's feature size.

Fourier's law for non-uniform conductivity systems is given by
\begin{align}
    \nabla \kappa(x, y) \nabla \Delta T(x, y) &= 0
\end{align}
% https://www.osti.gov/servlets/purl/1570911
Here $\Delta T$ is the local temperature deviation from the equilibrium of the system. At the nanoscale we need Boltzmann's equation. 
\begin{align}
    \Lambda \tilde{s} \cdot \nabla \Delta T(\phi) + \Delta T(\phi) &= \frac{1}{2\pi} \int_0^{2\pi} \Delta T(\phi') d\phi'
\end{align}
Phonon directions are parameterized by polar angle $\phi$ and $\Delta T(x, y, \phi)$ is a quantity commonly referred to as the phonon temperature. 
The main difference here is that we need to take into account the phonon mean free path $\Lambda \tilde{s}$.



\section{The Three Field Approach}

The three-field approach is a standard method in the field. Using a finite grid can create checkerboard patterns. The base field $\rho$ is filtered into field $\tilde{\rho}$. This filtered field is then typically thresholded into field $\bar{\rho}$. 

%\section{Solving the Boltzmann Transport Equation (BTE)}

%Challenging the Boltzmann Transport equation raises a number of challenges. For start, what is the control variable?

\section{Transmission Interpolation Method}

Solving the Boltzmann transport equation becomes especially challenging in settings with boundaries. The transmission interpolation model (TIM) is one approach to handle such systems. This approach relies on the introduction of an ficitious density $\rho(r)$ used to track some property of interest such as thermal conductivity. We have
\begin{align}
    \rho \in [0, 1]^N
\end{align}
where $N$ is the discretized size of the grid we consider. To model a boundary, we introduce a quantity known as the transmission coefficient $t$, defined as
\begin{align}
    t &= \frac{2 \rho(r^-) \rho(r^+)}{\rho(r-) + \rho(r+)}
\end{align}
We then create an imaginary interface. If $\rho(r^-) = \rho(r^+)$ then $t=1$ or $t=0$. Here $r^+$ and $r^-$ are the two sides of the boundary. %\textcolor{red}{TODO: Verify}

\subsection{Boundary Conditions}

We apply a temperature gradient on the direction we are interested in. We need to make our boundary condition translation invariant. We create a forward solver

\begin{align}
\Lambda \tilde{s} \cdot \nabla \Delta T(x,y, \phi) + \Delta T(x, y,\phi) &= \frac{1}{2\pi} \int_0^{2\pi} \Delta T(x, y, \phi') d\phi'
\end{align}
We transform this into a matrix linear system.
\begin{align}
    \sum_{p'q'} G^{p'q'}_{p q} \Delta T_{p' q'} &= P_{pq}
\end{align}
$P$ is the externally applied temperature gradient. We can write
\begin{align}
    G^{p' q'}_{[p] q} &= \delta_{q q'} \delta_{p p'} (1 + D_{p' q'} ) + \delta_{p p'} A_{p q q'} + \delta_{q q'} (B_{qp p'} - \frac{1}{M})
\end{align}
where $A$ is the scattering operator and $B$ is auxiliary tensor. We create a linear operator
\begin{align}
    L(X)_{p q} &= X_{p q}D_{p q} + \Sigma_{q'} A_{p q q'}X_{p q' } + \sum (B_{q p p'}-\frac{1}{M})X_{p'q}
\end{align}

This turns into a system of linear equations which can be solved with standard techniques. %We also need to develop a Jacobi precondition.

%\begin{align}
%    L(\Delta T, \bar{\rho}) &= P(\bar{\rho})
%\end{align}
%We also have to deal with the scaled thermal conductivity

%\begin{align}
%f_{\kappa}(\bar{\rho}) &= ...
%\end{align}

%Scipy's GMRES is used but needs some flattening to make everything work.

\subsection{Compute Gradients with Adjoint Method}

We use the adjoint method to compute the gradient of this system 
\begin{align}
    \frac{d \bar{\kappa}}{d \rho_p} &= \sum_{\mu c p} \left ( \frac{\partial \bar{\kappa}}{\partial \Delta \bar{T}_{\mu c}} \frac{\partial \Delta \bar{T}_{\mu c}}{\partial \bar{\rho}} \right) \frac{\partial \rho_{p'}}{\partial \rho_p} + \sum_{p'} \frac{\partial \bar{\kappa}^s}{\partial \bar{\rho}_{p'}} \frac{\partial \rho_{p'}}{\partial \rho_p}
\end{align}
% (\textcolor{red}{TODO: When is the adjoint method applicable?}
The solver is implemented with a combination of automatic differentiation and adjoint method gradient computation.

\section{Optimizing Thermal Transport}

Assume that $\rho(x, y)$ is the material density. Assuming Fourier's law, we can do the optimization
\begin{align}
    &\min_\rho g(\rho)\\
    &\textrm{s.t.}\quad \nabla \cdot \kappa(x, y) \nabla \Delta T(x, y) = 0
\end{align}
Gradients can be computed using the adjoint approach, which allows for constrained optimization.

\subsection{Optimizing for Heat Management}

The task here is to design a material with the desired effective thermal constraint.
\begin{align}
    &\min_\rho |\bar{\kappa}(\rho) - \tilde{\kappa} |_{Fro} \\
    &\textrm{s.t } 0 \leq \rho \leq 1 \\
    &\textrm{s.t } \sum_n \bar{\rho}_n N \geq \varphi  
\end{align}
The BTE solver gives us $\bar{\kappa}(\bar{\rho})$. This joint system can be used to optimize heat flows in nanostructures.
% (\textcolor{red}{TODO: Are there applications to chip design we should highlight})

The phonon mean free path is much more spread out than the electron free path. This means the BTE paradigm is very useful for thermoelectrics
\end{document}